કેટલીક વસ્તુઓના સમૂહને \emph{ગણ} કહે છે. તે વસ્તુઓ ગણના ઘટકો છે. જો $x$ ગણ $A$નો ઘટક હોય તો $x \in A$ લખીએ છીએ. ગણ \emph{વિસ્તારાત્મક} છે, એટલે તેના ઘટકો તેને સંપૂર્ણપણે નિર્ધારિત કરે છે. ઘટકોને સ્પષ્ટ રીતે \emph{યાદીમાં મૂકીને} અથવા તેમનામાં સામાન્ય હોય એવો ગુણધર્મ આપીને (\emph{અમૂર્તીકરણથી}) ગણ નિર્દિષ્ટ કરી શકાય છે. વિસ્તારાત્મકતાને લીધે ઘટકો કયા ક્રમે અથવા કઈ રીતે દર્શાવ્યા છે તે અગત્યનું નથી. $A$ અને $B$ એક જ ગણ છે ($A=B$) એમ સાબિત કરવા $A$નો દરેક ઘટક $B$માં અને $B$નો દરેક ઘટક $A$માં છે એમ બતાવવું પડે છે.

મહત્ત્વના ગણોમાં પ્રાકૃતિક ($\Nat$), પૂર્ણાંક ($\Int$), સંમેય ($\Rat$) અને વાસ્તવિક ($\Real$) સંખ્યાઓ તેમજ વસ્તુઓની \emph{શ્રેણીઓ} ($X^*$) અને અનંત \emph{અનુક્રમો} ($X^\omega$) આવે છે. $A$નો દરેક ઘટક $B$નો પણ ઘટક હોય તો $A$ એ $B$નો \emph{ઉપગણ} છે, $A\subseteq B$. $B$ના બધા ઉપગણોનો ગણ તેનો \emph{ઘાતગણ} $\Pow{B}$ છે. ગણોનો \emph{યોગગણ} $A\cup B$ અને \emph{છેદગણ} $A\cap B$ બનાવી શકાય છે. \emph{ક્રમિત જોડ} $\tuple{x,y}$માં $x$ અને $y$ એ જ ક્રમમાં આવે છે. $x=y$ ન હોય તો $\tuple{x,y}$ અને $\tuple{y,x}$ જુદી જોડો છે. $x\in A$ અને $y\in B$ હોય એવી બધી જોડો $\tuple{x,y}$નો ગણ $A$ અને $B$નો \emph{કાર્તેઝીય ગુણાકાર} $A\times B$ છે. $A\times A$ માટે $A^2$ લખીએ છીએ. દાખલા તરીકે $\Nat^2$ પ્રાકૃતિક સંખ્યાઓની જોડોનો ગણ છે.
